\section{Выводы}

В ходе работы я пришёл к следующим выводам:
\begin{itemize}
    \item \textbf{Эффективность суффиксного массива:} Суффиксный массив оказался подходящей структурой данных для задачи поиска множества образцов в заранее известном тексте. После однократного построения структуры можно быстро обрабатывать большое количество запросов.

    \item \textbf{Роль массива LCP:} Использование массива LCP позволило избежать повторного сравнения одних и тех же префиксов при бинарном поиске по суффиксному массиву. В обычной реализации поиск мог занимать \(O(m \log n)\), так как на каждом шаге бинарного поиска образец мог сравниваться с суффиксом с самого начала. Добавление LCP-оптимизации улучшило оценку поиска до \(O(m + \log n)\).

    \item \textbf{Использование RMQ:} Для быстрого получения LCP двух произвольных суффиксов была построена структура RMQ на основе Sparse Table. Это позволило отвечать на запросы минимума на отрезке массива LCP за \(O(1)\), что необходимо для LCP-ускоренного бинарного поиска.

    \item \textbf{Экспериментальное подтверждение сложности:} В ходе бенчмарка длина текста увеличивалась с \(10000\) до \(1280000\), то есть в \(128\) раз, при фиксированной длине образца \(m = 64\). При этом среднее время одного поиска выросло только с \(0.117873\) до \(0.156424\) микросекунды. Это показывает, что поиск не зависит от длины текста линейно и согласуется с оценкой \(O(m + \log n)\).

    \item \textbf{Стоимость предварительной обработки:} Построение суффиксного массива, массива LCP и RMQ требует дополнительного времени и памяти. Однако эта стоимость оправдана, так как все структуры строятся один раз для заранее известного текста, после чего используются для обработки всех последующих образцов.

    \item \textbf{Практические особенности реализации:} Помимо асимптотической сложности, важную роль играют детали реализации: сортировка найденных позиций для вывода в возрастающем порядке, буферизация вывода и аккуратная работа с индексами. Эти детали не меняют основной идеи алгоритма, но влияют на прохождение задачи в условиях ограничения времени.
\end{itemize}

Данный опыт закрепил понимание того, что при решении задач на строки важно учитывать не только саму структуру данных, но и вспомогательные алгоритмы вокруг неё. Суффиксный массив сам по себе позволяет выполнять поиск образцов, однако использование LCP и RMQ делает этот поиск теоретически и практически более эффективным. Также работа показала, что экспериментальные замеры полезны для проверки того, насколько реальное поведение программы соответствует заявленной асимптотической оценке.

\pagebreak